\begin{tikzpicture}

    % --- 1. VRSTICA ---

    % T1: id v Q0 in P0
    \node[table_node] (T1) {
        $\begin{array}{c|ccccc}
            & P_0 & P_1 & P_2 & P_3 & P_4 \\
            \hline 
            Q_0 & \text{id} & \text{id} & \text{id} & \text{id} & \text{id} \\
            Q_1 & \text{id} & & & & \\
            Q_2 & \text{id} & & & & \\
            Q_3 & \text{id} & & & & \\
            Q_4 & \text{id} & & & & \\
        \end{array}$
    };

    % T2: id + Q1 vrstica in P1 stolpec
    \node[table_node, right=2.2cm of T1] (T2) {
        $\begin{array}{c|ccccc}
            & P_0 & P_1 & P_2 & P_3 & P_4 \\
            \hline 
            Q_0 & \text{id} & \text{id} & \text{id} & \text{id} & \text{id} \\
            Q_1 & \text{id} & \rho & \rho^2 & \rho^3 & \rho^4 \\
            Q_2 & \text{id} & \rho^2 & & & \\
            Q_3 & \text{id} & \rho^3 & & & \\
            Q_4 & \text{id} & \rho^4 & & & \\
        \end{array}$
    };

    % --- 2. VRSTICA ---

    % T3: Dodan element Q2, P2
    \node[table_node, below=1.2cm of T1] (T3) {
        $\begin{array}{c|ccccc}
            & P_0 & P_1 & P_2 & P_3 & P_4 \\
            \hline 
            Q_0 & \text{id} & \text{id} & \text{id} & \text{id} & \text{id} \\
            Q_1 & \text{id} & \rho & \rho^2 & \rho^3 & \rho^4 \\
            Q_2 & \text{id} & \rho^2 & & & \\
            Q_3 & \text{id} & \rho^3 & \rho & & \\
            Q_4 & \text{id} & \rho^4 & & & \\
        \end{array}$
    };

    % T4: Končna tabela
    \node[table_node, right=2.2cm of T3] (T4) {
        $\begin{array}{c|ccccc}
            & P_0 & P_1 & P_2 & P_3 & P_4 \\
            \hline 
            Q_0 & \text{id} & \text{id} & \text{id} & \text{id} & \text{id} \\
            Q_1 & \text{id} & \rho & \rho^2 & \rho^3 & \rho^4 \\
            Q_2 & \text{id} & \rho^2 & \rho^4 & \rho & \rho^3 \\
            Q_3 & \text{id} & \rho^3 & \rho & \rho^4 & \rho^2 \\
            Q_4 & \text{id} & \rho^4 & \rho^3 & \rho^2 & \rho \\
        \end{array}$
    };

    % --- PUŠČICE ---

    % Ravna 1 -> 2
    \draw[main_arrow] (T1.east) -- (T2.west);

    % Zavita 2 -> 3 (iz spodnjega levega kota T2 v zgornji desni kot T3)
    % Uporabimo in/out kote za gladek prehod
    \draw[main_arrow] (T2.south west) to[out=225, in=45] (T3.north east);

    % Ravna 3 -> 4
    \draw[main_arrow] (T3.east) -- (T4.west);

\end{tikzpicture}